% abstract, choose between abstract and summary

Open source software is the foundation of modern computing, and almost every codebase
in use today depends on it. The way vulnerabilities are handled in this ecosystem
creates a structural problem. Under the Coordinated Vulnerability Disclosure model a
flaw is corrected before it is announced, and the fixing commit on GitHub is deliberately kept
quiet, carrying no reference to security. The correction is therefore public and
readable in the repository, while the information needed to recognise it as a security
fix is not.

This opens a window in which the projects that depend on that code are exposed: 
the fix cannot be installed until a new release is published, and nothing
notifies users until an advisory is issued. Meanwhile, the change itself describes the
vulnerability, showing which lines were wrong and how they were repaired, so that
anyone who reads it can work backwards to the flaw. 
Reading unfamiliar code and judging its security has always required expertise and time but 
general purpose language models may be lowering both barriers. 

Investigating this question requires a collection of repositories whose vulnerabilities
are known, and assembling one is far from simple. 
Public vulnerability records rarely
point to the commit that resolves them, and when they do the reference is often broken,
while the fixing commits cannot be found by searching either, if they are silent patch, then
there is no sign of what it fixes.

Earlier approaches applied deep learning to source code, while more recent ones 
turn language models into agents that explore a
repository by themselves, and some of these systems have already found real
vulnerabilities. The systems that succeed, however, rely on specialised infrastructure, 
while a model operating on its own performs considerably worse. This thesis examines the 
intermediate configuration, which reflects the resources realistically available to an ordinary user.

This thesis proposes an architecture based on AI agents that inspect a repository on
their own, decide whether it contains a vulnerability and classify it with the
corresponding CWE. Each agent runs inside a disposable container, with read-only access
to the code and no access to the web. Three configurations are compared, differing only
in what the agent is allowed to access: the code changed by a commit together with its
diff file, those same changes plus the whole repository, and the repository alone with no
indication that a fix exists, which reproduces the search for an unknown vulnerability.
The dataset is built from repositories that still contained the
vulnerability, one commit before the fix, and is restricted to recent disclosures so
that the answers cannot simply be retrived from training data.

As a secondary contribution, the thesis prepares the basis for a real-time extension. A
monitoring tool selects the repositories that are worth tracking, 
recording which of them end up with a published vulnerability and how long after
the selection. This is the component on which a continuous version of the analysis would
be built, moving it from vulnerabilities that are already known to code that has not yet been the
subject of any report.
